\PassOptionsToPackage{dvipsnames}{xcolor}

\documentclass[acmsmall,screen,review]{acmart}

\setcopyright{acmlicensed}
\copyrightyear{2018}
\acmYear{2018}
\acmDOI{XXXXXXX.XXXXXXX}
\acmConference[Conference acronym 'XX]{Make sure to enter the correct
  conference title from your rights confirmation emai}{June 03--05,
  2018}{Woodstock, NY}
\acmISBN{978-1-4503-XXXX-X/18/06}

\citestyle{acmauthoryear}

\bibliographystyle{ACM-Reference-Format}

\usepackage{hyperref}
\usepackage{anyfontsize} % kill warnings about font sizes
\usepackage{amsmath}     % Mathematics FTW!
\usepackage[capitalise]{cleveref}
%\usepackage{amsthm}      % Warning, I've only proven it correct :).
%\usepackage{amssymb}     % needed for \mathbb and mathcal
%\usepackage{bbm}
\usepackage[english]{babel}
%\usepackage{amsbsy}      % for \boldsymbol
%\usepackage{stmaryrd}   % provides some PLy symbols, like Scott Brackets
\usepackage{mathpartir}  % for \mathpar
\usepackage{color}
\usepackage{braket}      % for \set and \braket
\usepackage{mathtools}   % for \bigtimes
\usepackage{thmtools}    % to duplicate theorems in the appendix
%\usepackage{marvosym} % for deniarius (gradual label)
%\usepackage{wasysym}     % for \smiley
%\usepackage{balance}     % for the references
\usepackage{xspace}
%\usepackage{mathrsfs} % for mathscr font
\usepackage{proof}
% \usepackage{yhmath}
\usepackage{rotating}
\usepackage[Symbol]{upgreek} % upright greek letters
\usepackage{tikz-cd}
\usepackage{xcolor}
\usepackage{changepage}
\usepackage{centernot}
\usepackage{comment}
\usepackage{listings}
\usepackage{mdframed}
\usepackage{centernot}
\usepackage{mathtools}
\usepackage{multirow}
\usepackage{xparse}
\definecolor{darkblue}{rgb}{0.0, 0.0, 0.55}
\definecolor{darkgreen}{rgb}{0.0, 0.55, 0.13}
\definecolor{darkred}{rgb}{0.55, 0.0, 0.0}
\definecolor{darkviolet}{rgb}{0.58, 0.0, 0.83}
\definecolor{lightblue}{rgb}{0.68, 0.85, 0.9}
\usepackage{lstcoq}
% Literate Agda
% \usepackage{agda}
\usepackage[all]{xy}
\usepackage{stmaryrd}
\usepackage{subcaption,verbatim}
\usepackage[utf8x]{inputenc}
% \usepackage{agda-syntax}
\usepackage[outline]{contour}


\DeclareMathAlphabet{\mathbbm}{U}{bbm}{m}{n}

%% hack to get upright greek letters
\renewcommand{\Omega}{\Upomega}
\renewcommand{\Pi}{\Uppi}
\renewcommand{\Sigma}{\Upsigma}
\renewcommand{\Gamma}{\Upgamma}
\renewcommand{\Delta}{\Updelta}

\makeatletter
% \@ifpackageloaded{stix}{%
% }{%
%   \DeclareFontEncoding{LS2}{}{\noaccents@}
%   \DeclareFontSubstitution{LS2}{stix}{m}{n}
%   \DeclareSymbolFont{stix@largesymbols}{LS2}{stixex}{m}{n}
%   \SetSymbolFont{stix@largesymbols}{bold}{LS2}{stixex}{b}{n}
%   \DeclareMathDelimiter{\lBrace}{\mathopen} {stix@largesymbols}{"E8}%
%                                             {stix@largesymbols}{"0E}
%   \DeclareMathDelimiter{\rBrace}{\mathclose}{stix@largesymbols}{"E9}%
%                                             {stix@largesymbols}{"0F}
% }
\makeatother
\DeclareMathAlphabet{\mathcal}{OMS}{cmsy}{m}{n} % Reset mathcal font because ACM version sucks

% Comment form:  requires amssymb and color

\newcommand{\mynote}[2]
    {{\color{red} \fbox{\bfseries\sffamily\scriptsize#1}
    {\small$\blacktriangleright$\textsf{\emph{#2}}$\blacktriangleleft$}}~}

\definecolor{propcolor}{HTML}{3F7D31}
\definecolor{mypurple}{HTML}{5B069D}
\newcommand{\myproposal}[2]
    {{\color{propcolor} \fbox{\bfseries\sffamily\scriptsize#1}
    {\small$\blacktriangleright$\textsf{\emph{#2}}$\blacktriangleleft$}}~}

\newcommand{\nt}[1]{\mynote{NT}{#1}}
\newcommand{\lp}[1]{\mynote{LP}{#1}}
\newcommand{\jc}[1]{\mynote{JC}{#1}}
\newcommand{\tw}[1]{\mynote{TW}{#1}}
\newcommand{\pmp}[1]{\mynote{PMP}{#1}}

% Grey Box
\definecolor{lightgray}{gray}{0.90}
\newcommand{\Gbox}[1]{\colorbox{lightgray}{$#1$}}

% text
\newcommand{\ie}{\emph{i.e.,}\xspace}
\newcommand{\etal}{\emph{
    et al.}\xspace}
\newcommand{\eg}{\emph{e.g.,}\xspace}

\newcommand{\nat}{\mathbb{N}}
\newcommand{\natP}{\mathbb{N}_\mathbb{P}}
\newcommand{\bool}{\mathsf{Bool}}

\newcommand{\slogan}[1]{\begin{center}``\emph{#1}''\end{center}}

\newcommand{\agdahtml}{https://htmlpreview.github.io/?https://github.com/CoqHott/logrel-mltt/blob/impredicativity-cast-compute-refl/html/}
\newcommand{\refAgdaRoot}[1]{{\small{\href{\agdahtml/#1.html}{\emph{[#1.agda]}}}}}
\newcommand{\refAgdaDef}[1]{{\small{\href{\agdahtml/Definition.#1.html}{\emph{[#1.agda]}}}}}
\newcommand{\refAgda}[2]{{\small{\href{\agdahtml/Definition.#1.#2.html}{\emph{[#2.agda]}}}}}
\newcommand{\refCoq}[1]{{\color{darkblue}\emph{[#1]}}}
% \newcommand{\vrefCoq}[2]{\href{https://github.com/TheoWinterhalter/template-coq/tree/rewrite-rules/pcuic/theories/#1}{\emph{[#2]}}}
\newcommand{\vrefCoq}[2]{\href{https://github.com/TheoWinterhalter/template-coq/tree/popl21-artifact/pcuic/theories/#1}{\emph{[#2]}}}


% \newcommand{\Coq}{${\mathrm{Coq}}$\xspace}
\newcommand{\MLTT}{\ensuremath{\mathsf{MLTT}}\xspace}
\newcommand{\sMLTT}{\ensuremath{\mathsf{sMLTT}}\xspace}
\newcommand{\ETT}{\ensuremath{\mathsf{ETT}}\xspace}
\newcommand{\HTS}{\ensuremath{\mathsf{HTS}}\xspace}
\newcommand{\SetoidTT}{\ensuremath{\mathsf{TT}^\mathrm{obs}}\xspace}
\newcommand{\SetoidCC}{\ensuremath{\mathsf{CC}^\mathrm{obs}}\xspace}
\newcommand{\SetoidCIC}{\ensuremath{\mathsf{CIC}^\mathrm{obs}}\xspace}
\newcommand{\CIC}{\ensuremath{\mathsf{CIC}}\xspace}

\newcommand{\UIP}{${\mathrm{UIP}}$\xspace}

\newcommand{\bbfamily}{\fontencoding{U}\fontfamily{bbold}\selectfont}
% \DeclareMathAlphabet{\mathbbold}{U}{bbold}{m}{n}
\newcommand{\Two}{\omega}
%% \newcommand{\Two}{{\color{white}\contour{blga}}
\newcommand{\independent}{\emptyset}

\newcommand{\prop}{\prim{Prop}}
\newcommand{\Univ}{\oblset{Univ}}
\newcommand{\hProp}{\oblset{hProp}}
\def\Coq{\textsc{Coq}\xspace}
\def\Agda{\textsc{Agda}\xspace}
\def\Lean{\textsc{Lean}\xspace}
\def\Idris{\textsc{Idris}\xspace}
% \newcommand{\Agda}{${\mathrm{Agda}}$\xspace}
% \newcommand{\Lean}{${\mathrm{Lean}}$\xspace}
% \newcommand{\Cedille}{${\mathrm{Cedille}}$\xspace}
\newcommand{\CICU}{${\mathrm{CIC}}_u$\xspace}
\newcommand{\RTT}{${\mathrm{RTT}}$\xspace}
\def\MetaCoq{\textsc{MetaCoq}\xspace}

\newcommand{\subst}[3][x]{#2[{#1}:={#3}]}

\newcommand\dashPar{\vdash}%_{p}}

\newcommand{\sizeName}{\mathtt{size}}
\newcommand{\size}[1]{\sizeName(#1)}

\newcommand{\wfctx}[2][]{#1 \vdash #2}
\newcommand{\eqctx}[2]{\vdash #1 \equiv #2}
\newcommand{\tytm}[4][]{#1 #2 \vdash #3 : #4}
\newcommand{\tytmPara}[4][]{#1 #2 \dashPar #3 : #4}
\newcommand{\tm}[2]{#1 : #2}
\newcommand{\tytms}{\tytm}
\newcommand{\eqtm}[5][]{#1 #2 \vdash #3 \equiv #4 : #5}
\newcommand{\nered}[4]{#1 \vdash #2 \Rightarrow_{ne} #3 : #4}
\newcommand{\neconv}[4]{#1 \vdash #2 \cong_{ne} #3 : #4}
\newcommand{\neconvRed}[4]{#1 \vdash #2 \cong_{ne}^\downarrow #3 : #4}
\newcommand{\whnfconv}[4]{#1 \vdash #2 \cong #3 : #4}
\newcommand{\whnfconvRed}[4]{#1 \vdash #2 \cong^\downarrow #3 : #4}
\newcommand{\red}[5][]{#1 #2 \vdash #3 \Rightarrow #4 : #5}
\newcommand{\redT}[5][]{#1 #2 \vdash #3 \Rightarrow^* #4 : #5}
\newcommand{\eqtmmultiline}[6][]{#1 #2 \vdash \begin{array}{l} #3
                                                \equiv \\
                                                #4 \end{array} : #5 : #6}
\newcommand{\redmultiline}[5][]{#1 #2 \vdash \begin{array}{l} #3 \Rightarrow \\ #4 \end{array} : #5}
\newcommand{\redmultilineExt}[6][]{#1 #2 \vdash \begin{array}{l} #3 \Rightarrow \\ #4 \\ #5 \end{array} : #6}

\newcommand{\eqproj}[3]{#1_\epsilon^{#2}~#3}
\newcommand{\propext}{\sProp{}_\mathrm{ext}}
\newcommand{\funext}{\Pi_\mathrm{ext}}

\newcommand{\tytyR}[4]{#2 \Vdash_{#1} #3 : #4}
\newcommand{\tytmR}[5]{#2 \Vdash_{#1} #3 : #4 : #5}
\newcommand{\ctxV}[1]{\Vdash^v #1}
\newcommand{\tytyV}[4]{#2 \Vdash^v_{#1} #3 : #4}
\newcommand{\tytyS}[3]{#1 \Vdash^s #2 : #3}
\newcommand{\tytmV}[5]{#2 \Vdash^v_{#1} #3 : #4 : #5}
\newcommand{\nattmR}[2]{#1 \Vdash_{\mathbb{N}\mathsf{t}}\ {#2}}
\newcommand{\nateqR}[3]{#1 \Vdash_{\mathbb{N}\mathsf{t}=}\ {#2} \equiv {#3}}
\newcommand{\quotmR}[2]{#1 \Vdash_Q\ {#2}}
\newcommand{\quoeqR}[3]{#1 \Vdash_Q\ {#2} \equiv {#3}}
\newcommand{\idtmR}[2]{#1 \Vdash_\mathrm{Id}\ {#2}}
\newcommand{\ideqR}[3]{#1 \Vdash_\mathrm{Id}\ {#2} \equiv {#3}}
\newcommand{\eqtyR}[5]{#2 \Vdash_{#1} #3 \equiv #4 : #5}
\newcommand{\eqtyS}[4]{#1 \Vdash^s #2 \equiv #3 : #4}
\newcommand{\eqtmR}[6]{#2 \Vdash_{#1} #3 \equiv #4 : #5 : #6}
\newcommand{\eqtmV}[6]{#2 \Vdash^v_{#1} #3 \equiv #4 : #5 : #6}
\newcommand{\wk}[3]{{#1}}% : {#2} \le {#3}}
\newcommand{\wkrho}{(\rho : \Delta \sqsubseteq \Gamma)}
\newcommand{\wkrhoshort}{\wkrho}

\newcommand{\metatm}[2]{{#1} :: {#2}}
\newcommand{\metaforall}[2]{\forall {#1},\ {#2}}
\newcommand{\metaimply}[2]{{#1}\ \to\ {#2}}

\newcommand{\emptyctx}{\bullet}
\newcommand{\extctx}[3][x]{#2 , \tm{#1}{#3}}
\newcommand{\inctx}[4][x]{\tmannotated{#1}{#3}{#4}{} \in {#2}}

\newcommand{\Type}{\mathsf{Type}}
\newcommand{\SProp}{\mathsf{SProp}}
\newcommand{\Prop}{\mathsf{Prop}}

\newcommand{\forded}[1]{#1 _\mathbb{F}}
\newcommand{\eqpi}[2][]{#2 _\epsilon^{#1}}
\newcommand{\appi}[1]{#1 _{ap}}
\newcommand{\Depsum}[3][x]{\Sigma (\tm{#1}{#2}).\, {#3}}
\newcommand{\Depfun}[3][x]{\Pi (\tm{#1}{#2}).\, {#3}}
\newcommand{\Fun}[2]{{#1} \to {#2}}
\newcommand{\Sig}[3][x]{\Sigma (\tm{#1}{#2}).\, {#3}}
\newcommand{\Prod}[2]{{#1} \land {#2}}
\newcommand{\Exists}[3][x]{(\tm{#1}{#2}) \ \& \ {#3}}
\newcommand{\varExists}[3][x]{\exists (\tm{#1}{#2})\ .\ {#3}}
\newcommand{\ExistsUnique}[3][x]{\exists! (\tm{#1}{#2})\ .\ {#3}}
\newcommand{\Path}[3][]{#2 \sim_{#1} #3}
\newcommand{\eqInd}[3]{\mathrm{eq}_{#1} \, #2 \, #3}
\newcommand{\Empty}{\bot}
\newcommand{\botP}{\bot_{\mathbb{P}}}
\newcommand{\Unit}{\top}
\newcommand{\Quo}[2]{{#1}/{#2}}
\newcommand{\Id}[3]{\metaop{Id}(#1,#2,#3)}
\newcommand{\Boxt}[1]{\square{#1}}
\newcommand{\Squash}[1]{\| #1 \|}
\newcommand{\CubicalPath}{\metaop{Path}}
\newcommand{\List}[1]{\mathtt{list} \, #1}

\newcommand\metaop[1]{{\color{RoyalBlue}{\mathrm{#1}}}}

\newcommand{\piRel}[2]{\mathcal{R} ( #1 , #2)}
\newcommand{\univRel}[2]{\mathcal{A} ( #1 , #2)}

\newcommand{\lstop}[1]{\textrm{\coqeN{#1}}\xspace}
\newcommand{\lstops}[1]{\textrm{\lstinline[language=Coq, basicstyle=\small,escapechar=*]{#1}}\xspace}

\newcommand{\emptyrec}[2]{\mathtt{\color{darkblue} \bot\mathrm{-} elim}~{#1}~{#2}}
\newcommand{\unit}{*}
\newcommand{\zero}{{\color{darkblue} 0}}
\newcommand{\sucName}{\mathtt{\color{darkblue} S}}
\newcommand{\suc}[1]{\sucName\ #1}
\newcommand{\natrec}[4]{\mathtt{\color{darkblue} \mathbb{N}\mathrm{-} elim}~{#1}~{#2}~{#3}~{#4}}
\newcommand{\lam}[3][x]{\lambda (\tm{#1}{#2}).\, #3}
\newcommand{\pair}[2]{\langle #1 , #2 \rangle }
\newcommand{\fst}[1]{\metaop{fst}({#1})}
\newcommand{\snd}[1]{\metaop{snd}({#1})}
% \newcommand{\refl}[2]{\lstop{refl}~#1}
\newcommand{\refl}[2]{\mathtt{\color{darkblue} refl}~#1}
\newcommand{\reflI}[2]{\mathtt{\color{darkblue} refl}~#1}
\newcommand{\sym}[1]{{#1}^{-1}}
\newcommand{\transitivity}[2]{{#1}\cdot{#2}}
\newcommand{\ap}[2]{\lstop{ap}~{#1}~{#2}}
\newcommand{\castName}{\lstop{cast}}
% \newcommand{\cast}[4]{{\castName~#1~#2~#3~#4}}
\newcommand{\cast}[4]{{\mathtt{\color{darkblue} cast}~#1~#2~#3~#4}}
\newcommand{\castI}[4]{{\mathtt{\color{darkblue} cast}~#1~#2~#3~#4}}
\newcommand{\transportName}{\lstop{transp}}
% \newcommand{\transport}[6]{\transportName~#1~#2~#3~#4~#5~#6}
\newcommand{\transport}[6]{\mathtt{\color{darkblue} transport}~#1~#2~#3~#4~#5~#6}
\newcommand{\transportI}[6]{\mathtt{\color{darkblue} transport}~#1~#2~#3~#4~#5~#6}
\newcommand{\castreflName}{\metaop{castrefl}}
\newcommand{\castrefl}[2]{\castreflName(#1,#2)}
\newcommand{\quo}[1]{\pi(#1)}
\newcommand{\quorec}[4]{\metaop{Q\mathrm{-}elim}(#1,#2,#3,#4)}
\newcommand{\id}[1]{\metaop{Idrefl}(#1)}
\newcommand{\idpath}[1]{\metaop{Idpath}(#1)}
\newcommand{\J}[6]{\metaop{J}(#1,#2,#3,#4,#5,#6)}
\newcommand{\eqJ}[5]{\metaop{eq_J}(#1,#2,#3,#4,#5)}
\newcommand{\boxt}[1]{\diamond{#1}}
\newcommand{\unboxt}[1]{\metaop{\Box\mathrm{-elim}}(#1)}
\newcommand{\squash}[1]{| #1 |}
\newcommand{\unsquash}[3]{\metaop{\mathrm{S}\mathrm{-elim}}(#1,#2,#3)}
\newcommand{\nil}{\mathtt{nil}}
\newcommand{\cons}[2]{\mathtt{cons} \, #1 \, #2}

\newcommand\head[1]{\mathop{\text{hd}}#1}
\newcommand\whnf[1]{\mathop{\text{whnf}}#1}
\newcommand\neutral[1]{\mathop{\text{neutral}} #1}
\newcommand\postype[1]{\mathop{\text{posType}}#1}



\newcommand{\tmannotated}[3]{#1 : #2 : #3}%{#1 :^{#3} #2}
\newcommand{\extctxannotated}[4][x]{#2 , \tmannotated{#1}{#3}{#4}{}}
\newcommand{\tytmannotated}[5][]{#1 #2 \vdash \tmannotated{#3}{#4}{#5}}
\newcommand{\tytmParaannotated}[5][]{#1 #2 \dashPar \tmannotated{#3}{#4}{#5}}
\newcommand{\eqtmannotated}[6][]{#1 #2 \vdash #3 \equiv #4 : #5 : #6}
\newcommand{\algoeqtmannotated}[6][]{#1 #2 \vdash #3 \cong^\downarrow #4 : #5 : #6}
\newcommand{\eqtmParaannotated}[6][]{#1 #2 \dashPar #3 \equiv #4 : #5 : #6}
\newcommand{\redannotated}[5][]{#1 #2 \vdash #3 \Rightarrow #4 : #5}
\newcommand{\Depfunannotated}[6][x]{\Pi_{}^{#5, #6} (\tm{#1}{#2}).\, {#3}}
\newcommand{\Depsumannotated}[6][x]{\Sigma_{#4,#5}^{#6} (\tm{#1}{#2}).\, {#3}}
\newcommand{\Existsannotated}[5][x]{(\tm{#1}{#2}) \ \& \ {#3}}

\newcommand{\Con}{\mathsf{Con}}
\newcommand{\cwfty}[1][]{\mathsf{Ty}_{#1}}
\newcommand{\cwftm}[1][]{\mathsf{tm}_{#1}}
\newcommand{\cwfprop}[1][]{\mathsf{Pr}_{#1}}
\newcommand{\cwfpr}[1][]{\mathsf{pr}_{#1}}

\newcommand{\Setoid}{\mathsf{Setoid}}
\newcommand{\metasProp}{\mathsf{Prop}}
\newcommand{\metanat}{\boldsymbol{\omega}}
\newcommand{\metazero}{0}
\newcommand{\metasuc}[1]{\mathsf{S}\ #1}
\newcommand{\metaempty}{\bot}
\newcommand{\metaunit}{\top}
\newcommand{\metarefl}{\mathsf{refl}}

\renewcommand{\ne}{N}

\newcommand{\Context}{\AgdaData{Context}}
\newcommand{\Term}{\AgdaData{Term}}
\newcommand{\Sort}{\AgdaData{Sort}}

\newcommand{\VU}{{\textcolor{AgdaInductiveConstructor} \Vdash}_{\AgdaCon{U}}}
\newcommand{\VOmega}{{\textcolor{AgdaInductiveConstructor} \Vdash}_{\textcolor{AgdaInductiveConstructor} \Omega}}
\newcommand{\Vempty}{{\textcolor{AgdaInductiveConstructor} \Vdash}_{\textcolor{AgdaInductiveConstructor} \bot}}
\newcommand{\Vpi}{{\textcolor{AgdaInductiveConstructor} \Vdash}_{\textcolor{AgdaInductiveConstructor} \Pi}}
\newcommand{\Vforall}{{\textcolor{AgdaInductiveConstructor} \Vdash}_{\textcolor{AgdaInductiveConstructor} \forall}}
\newcommand{\Vexists}{{\textcolor{AgdaInductiveConstructor} \Vdash}_{\textcolor{AgdaInductiveConstructor} \&}}
\newcommand{\Vnat}{{\textcolor{AgdaInductiveConstructor} \Vdash}_{\AgdaCon{N}}}
\newcommand{\Vne}{{\textcolor{AgdaInductiveConstructor} \Vdash}_{\AgdaCon{ne}}}
\newcommand{\Vemb}{{\textcolor{AgdaInductiveConstructor} \Vdash}_{\AgdaCon{emb}}}
\newcommand{\VX}{{\textcolor{AgdaInductiveConstructor} \Vdash}_{\AgdaCon{X}}}

\newcommand{\RU}{\AgdaCon{R}_{\AgdaCon{U}}}
\newcommand{\ROmega}{\AgdaCon{R}_{\textcolor{AgdaInductiveConstructor} \Omega}}
\newcommand{\Rempty}{\AgdaCon{R}_{\textcolor{AgdaInductiveConstructor} \bot}}
\newcommand{\Rpi}{\AgdaCon{R}_{\textcolor{AgdaInductiveConstructor} \Pi}}
\newcommand{\Rforall}{\AgdaCon{R}_{\textcolor{AgdaInductiveConstructor} \forall}}
\newcommand{\Rexists}{\AgdaCon{R}_{\textcolor{AgdaInductiveConstructor} \exists}}
\newcommand{\Rnat}{\AgdaCon{R}_{\AgdaCon{N}}}
\newcommand{\Rne}{\AgdaCon{R}_{\AgdaCon{ne}}}
\newcommand{\Remb}{\AgdaCon{R}_{\AgdaCon{emb}}}

\newcommand{\sV}{\mathbf{V}}
\newcommand{\sU}{\mathbf{U}}
\newcommand{\Upred}[1][i]{\mathsf{code}_{#1}}
\newcommand{\sType}{\mathbf{s}}
\newcommand{\ssProp}{\Omega}
\newcommand{\sFun}[2]{#1 \to_s #2}
\newcommand{\sEq}[2]{#1 = #2}
\newcommand{\el}{\mathsf{el}}
\newcommand{\eleq}{\mathsf{el\text{-}eq}}
\newcommand{\val}{\mathsf{val}}
\newcommand{\valeq}{\mathsf{val\text{-}eq}}
\newcommand{\codeemb}{\mathsf{c}_\mathsf{emb}}
\newcommand{\codenat}{\mathsf{c}_\mathsf{N}}
\newcommand{\codepi}{\mathsf{c}_{\Pi}}
\newcommand{\codepiuu}{\mathsf{c}_{\Pi\mathsf{UU}}}
\newcommand{\codepisu}{\mathsf{c}_{\Pi\Omega\mathsf{U}}}
\newcommand{\codepius}{\mathsf{c}_{\Pi\mathsf{U}\Omega}}
\newcommand{\codepiss}{\mathsf{c}_{\Pi\Omega\Omega}}
\newcommand{\codeex}{\mathsf{c}_{\exists}}
\newcommand{\codeU}{\mathsf{c}_\mathsf{U}}
\newcommand{\codesProp}{\mathsf{c}_{\sProp}}
\newcommand{\codeInd}{\mathsf{c}_\mathsf{ind}}
\newcommand{\codeempty}{\mathsf{c}_\bot}
\newcommand{\codequo}{\mathsf{c}_Q}

\newcommand{\mctx}[1]{\llbracket\ #1\ \rrbracket}
\newcommand{\mty}[2][\Gamma]{\llbracket {#2} \rrbracket_{#1}}
\newcommand{\mtm}[3][\Gamma]{[\tm{#2}{#3}]_{#1}}

\newcommand{\bnfis}{::=}
\newcommand{\bnfin}{\in}
\newcommand{\sep}{|}
\newcommand{\qsep}{\ | \ }

% Automatically insert parentheses
\newcommand{\maybeparens}[2]{%
  \ifthenelse{\isempty{#2}}{#1}{%  if argument is empty
  \ifthenelse{\equal{#2}{(}}{#1#2}{%  or (, then do not parenthesize
  #1(#2)}}} % otherwise, parenthesize
\newcommand{\tdom}{\mathrm{dom}}
\newcommand{\dom}{\maybeparens{\tdom}}
\newcommand{\tPV}{\mathrm{PV}}
\newcommand{\PV}{\maybeparens{\tPV}}
\newcommand{\tFV}{\mathrm{FV}}
\newcommand{\FV}{\maybeparens{\tFV}}

\newcommand{\joinsub}{\uplus}%{\|}
\newcommand{\Joinsub}{\biguplus}%{\|}

\newcommand{\mathsc}[1]{\ensuremath{\textsc{#1}}}
\newcommand{\rulename}[1]{{\footnotesize\mathsc{#1}}}
\newcommand{\nru}[3]{\dfrac{\begin{array}{@{}c@{}}#1\end{array}}{\begin{array}{@{}c@{}}#2\end{array}}\ifthenelse{\isempty{#3}}{}{\ \rulename{#3}}}
\newcommand{\ru}[2]{\nru{#1}{#2}{}}

\newcommand{\myref}[2]{\hyperref[#2]{#1~\ref*{#2}}}
% Referencing figures
\newcommand{\figref}[1]{\myref{Fig.}{fig:#1}}
\newcommand{\secref}[1]{\myref{Sect.}{sec:#1}}
\newcommand{\defref}[1]{\myref{Definition}{def:#1}}
\newcommand{\lemref}[1]{\myref{Lemma}{lem:#1}}
\newcommand{\thmref}[1]{\myref{Theorem}{thm:#1}}
\newcommand{\eqnref}[1]{\hyperref[#1]{(\ref{#1})}}
\newcommand{\exref}[1]{\myref{Example}{ex:#1}}



% Agda highlighting macros
%% \newcommand{\data}[1]{{\AgdaDatatype{#1}}}
%% \newcommand{\record}[1]{{\AgdaRecord{#1}}}
%% \newcommand{\func}[1]{{\AgdaFunction{#1}}}
%% \newcommand{\prim}[1]{{\AgdaPrimitive{#1}}}
%% \newcommand{\primty}[1]{{\AgdaPrimitiveType{#1}}}
%% \newcommand{\sset}[1]{{\prim{Set}_{#1}}}
%% \newcommand{\var}[1]{{\AgdaBound{#1}}}
%% \newcommand{\keyw}[1]{{\AgdaKeyword{#1}}}
%% \newcommand{\con}[1]{{\AgdaInductiveConstructor{#1}}}
%% \newcommand{\field}[1]{{\AgdaField{#1}}}
%% \newcommand{\proj}[1]{\field{.#1}}
%% \newcommand{\pragma}[2]{\texttt{\{-\# } \AgdaKeyword{\texttt{#1}}\ #2 \texttt{ \#-\}}}

%% \newcommand{\Id}[1]{\mathrel{{\AgdaPrimitiveType{\equiv}}_{#1}}}
%% \newcommand{\Leq}{\mathrel{\AgdaDatatype{\leq}}}
%% \newcommand{\Simeq}{\mathrel{\AgdaRecord{\simeq}}}
%% \newcommand{\sumtype}{\mathrel{\AgdaDatatype{\uplus}}}
%% \newcommand{\prodtype}{\mathrel{\AgdaDatatype{\times}}}

%% \newcommand{\pat}[1]{\AgdaKeyword{?}#1}

%% \newcommand{\Nat}{{\color{AgdaDatatype}\mathbb{N}}}
%% \newcommand{\myNat}{\func{N}}
%% \newcommand{\vect}[2]{\func{Vec} \ #1 \ #2}
%% \newcommand{\myZero}{\func{zero}}
%% \newcommand{\Sone}{\func{$S^1$}}
%% \newcommand{\base}{\func{base}}
%% \newcommand{\Loop}{\func{loop}}
%% \newcommand{\mySucc}{\func{suc}}
%% \newcommand{\myNatRec}{\func{elim}}%_{\myNat}}
%% \newcommand{\myCatch}{\func{catch}}
%% \newcommand{\myRaise}{{\func{raise}}}
%% \newcommand{\Exc}{\AgdaPostulate{$\mathbb{E}$}}
%% \newcommand{\Int}{\data{\mathbb{Z}}}
%% \newcommand{\toptype}{{\color{AgdaDatatype}\top}}
%% \newcommand{\bottomtype}{\data{\bot}}
%% \newcommand{\sigmatype}{\data{\Sigma}}
%% \newcommand{\Bool}{{\color{AgdaDatatype}\mathbb{B}}}
%% \newcommand{\myMax}{\AgdaFun{max}}

%% \newcommand{\Fin}{\data{Fin}}
%% \newcommand{\fzero}{\con{fzero}}
%% \newcommand{\fsuc}{\con{fsuc}}
%% \newcommand{\lezero}{\con{lz}}
%% \newcommand{\lesuc}{\con{ls}}
%% \newcommand{\refl}{\con{refl}}
%% \newcommand{\reflbar}{\overline{\con{refl}}}
%% \newcommand{\zero}{\con{zero}}
%% \newcommand{\suc}{\con{suc}}
%% \newcommand{\true}{\con{true}}
%% \newcommand{\false}{\con{false}}
%% \newcommand{\J}{\prim{J}}
%% \newcommand{\K}{\prim{K}}

%% \newcommand{\zero}{{\color{AgdaDatatype}\mathsf{0}}}
%% \newcommand{\succ}{{\color{AgdaDatatype}\mathsf{S}}}
%% \newcommand{\ty}[1][]{{\color{AgdaDatatype}\square}_{#1}}
%% \newcommand{\natind}{\Nat_{\mathsf{ind}}}
%% \newcommand{\natind}{\myNatRec}

%% \newcommand{\rwRule}[4]{#1 \ctxrw \cmd{#2}{#3} \rwred #4}

%% \newcommand{\cmd}[2]{\langle #1 \mid #2 \rangle}
%% \newcommand{\cmdEval}[2]{\mathrm{eval}\, ( #1 \mid #2 )}
%% \newcommand{\zip}[2]{\mathsf{zip}(#1,\ #2)}
%% \newcommand{\zip}[2]{\mathsf{zip}(#1 \mid #2)}
%% \newcommand{\ctxrw}{\rightthreetimes}
%% \newcommand{\ctxrw}{\ltimes}
%% \newcommand{\confluent}{\downdownarrows}
%% \newcommand{\force}[1]{{\Downarrow}(#1)}
%% \newcommand{\patt}[1]{\AgdaKeyword{#1}}
%% \newcommand{\patt}[1]{{#1}}
%% \newcommand{\pattinst}[3]{#2 [ #3 ]_{#1}}
%% \newcommand{\force}[1]{{\patt{\Downarrow}}(#1)}
%% \newcommand{\erp}[1]{#1^{-}}
%% \newcommand\cored{\leftsquigarrow}
%% \newcommand{\rectype}[1]{\mathfrak{#1}}
%% \newcommand{\recmk}[2]{\lBrace{#2}\rBrace}
%% \newcommand{\recproj}[3]{{#2}._{#1} #3}
%% \newcommand{\stkproj}[2]{{\diamond}._{#1} #2}
%% \newcommand{\stkproj}[2]{{\diamond}.#2}
%% \newcommand{\stkapply}[1]{\diamond\ #1}
%% \newcommand{\inst}{\Rightarrow}
%% \newcommand{\patstk}{\mathbb{S}}
%% \newcommand{\patframe}{\mathbb{F}}

%% \newcommand{\reds}{\red^\star}
%% \newcommand{\coreds}{\cored^\star}

%% \newcommand\pos[2]{{% we make the whole thing an ordinary symbol
%%   \left.\kern-\nulldelimiterspace % automatically resize the bar with \right
%%   #1 % the function
%%   % \vphantom{\big|} % pretend it's a little taller at normal size
%%   \right|_{#2} % this is the delimiter
%%   }}


%% \newcommand\optRedName{\rho}
%% \newcommand\optRed[1]{\optRedName(#1)}
%% \def\paraRed{\Rrightarrow}
%% %\def\red{\rightarrow}
%% \newcommand\closure[1]{#1^*}
%% \def\Cred{\closure{\red}}

%% \newcommand\mgu[2]{\textit{MGU}(#1,#2)}

%% \newcommand{\RewContext}{\mathcal{R}}
%% \newcommand{\RewOrder}{\succ}

%% \newcommand\whnfName[1]{\mathtt{whnf}_{#1}}
%% \newcommand\whnf[2]{\whnfName{#1}(#2)}


%%%%%%%%%%%%%%%%%%%%%%%%%%%%%%%%%%%%%%%%%%%%%%%%%%%%%%%%%%%%%%%%%%%%%%%%%%%%%%%%%%%%%%%%%%%%%%%%%%%%%%%%%%%%%%%%%%%%%
%Reference to inference rules, inspired by https://tex.stackexchange.com/questions/340788/cross-referencing-inference-rules

\makeatletter
\let\originferrule\inferrule
\DeclareDocumentCommand \inferrule { s O {} m m }{%
	\IfBooleanTF{#1}%
	{%
		\mpr@inferstar[#2]{#3}{#4}%
	}{%
		\mpr@inferrule[#2]{#3}{#4}%
	}%
  % Modification of the SO code: we reuse the main label
	\IfValueT{#2}%
	{%
		\my@name@inferrule{#2}%
	}%
}
\NewDocumentCommand \my@name@inferrule { m }{%
	\def\@currentlabelname{\textsc{#1}}%
}

% \DeclareDocumentCommand \redrule {o m m}{%
% 	\IfValueTF{#1}
% 	{\textsc{#1}: \; #2 \redCCIC #3 \my@name@inferrule{#1}}
% 	{#2 \redCCIC #3}
% }

\newcommand*{\ilabel}[1]{%
  % \edef\@currentlabel{#1}% Set target label
  \phantomsection% Correct hyper reference link
  \label{#1}% Print and store label
}
\makeatother


\usepackage{ifthen}

\newboolean{appendix}
\setboolean{appendix}{false} % {true}
\newcommand{\ifAppendix}[2][]{\ifthenelse{\boolean{appendix}}{#2}{#1}}

\keywords{type theory, dependent types, rewriting theory, confluence, termination}

% Document starts
\begin{document}
% Title portion
\title{Weak Choice Principles for Proof-Relevant Setoids}
% \subtitle{A Type Theory with Computational Assumptions}

\author{Loïc Pujet}
\affiliation{%
  \institution{University of Stockholm}
  \city{Stockholm}
  \country{Sweden}
}

\newcommand{\shepherd}[1]
   {{#1}}
    % {{\color{red} #1}}

%
% The code below should be generated by the tool at
% http://dl.acm.org/ccs.cfm
% Please copy and paste the code instead of the example below.
%
\begin{CCSXML}
<ccs2012>
<concept>
<concept_id>10003752.10003790.10011740</concept_id>
<concept_desc>Theory of computation~Type theory</concept_desc>
<concept_significance>500</concept_significance>
</concept>
</ccs2012>
\end{CCSXML}

\ccsdesc[500]{Theory of computation~Type theory}
%
% End generated code
%

% \keywords{type theory, proof assistants, rewriting theory}


\begin{abstract}
Lorem ipsum dolor sit amet consectetuer adipiscing elit
\end{abstract}

\maketitle

\section{Introduction}

%
%% When it comes to writing formal proofs in dependent type theory, any seasoned
%% user of interactive proof assistants will tell you that picking the right
%% definition can often be more challenging than actually working out the proof.
%
Writing formal proofs in dependent type theory is a delicate art, even for the most
seasoned users of interactive proof assistants: the challenge often lies
not so much in working out the proof itself, but in finding the right
definition to build upon.
%
This is especially true in \emph{intensional} type theory~\citep{itt}, where the
equality judgment only identifies those definitions that are
\( \beta \eta \)-convertible.
%
Even though the intensional character of the judgmental equality can seem overly
restrictive compared to the extensional nature of the mathematical equality,
%
it is in fact a necessary condition for the decidability of the typing relation,
which is important for both foundational and practical reasons.
%
%% This is largely due to a mismatch between the \emph{extensional} aspects of
%% mathematics, where objects are characterised by their observational
%% behavior, and the \emph{intensional} nature of ITT, where objects are little
%% more than lambda-terms up to \( \beta \eta \)-equality.

The intensional nature of the equality judgment serves an important purpose,
but the same cannot be said of the \emph{propositional} equality,
which happens to be just as coarse.
%
As a result, users commonly postulate extensionality axioms for the
propositional equality, such as the principle of function
extensionality, which states that two functions are equal if their images are
equal, the principle of proposition extensionality, which states that
logically equivalent propositions are equal, the principle of uniqueness
of identity proofs (UIP), or even the univalence axiom.
%
Unfortunately, axioms tend to break the canonicity of the system: function
extensionality will result in closed terms of type \( \mathbb{N} \) that do not
evaluate to an actual integer.

An alternative strategy is to work with types that are equipped with an
equivalence relation, which are known as \emph{setoids}.
%
Doing so allows us to specify the meaning of equality for every type, but
it comes at the cost of having to manually prove that every function preserves
the setoid equality.
%
This should be the work of a compiler, not of the user!
%
What we really want is a way to interpret type theory with
extensionality principles into plain intensional type theory.
%
The first such interpretation was given by \citet{hofmann_setoids} with his setoid
model, which interprets all types as setoids and validates the principles of
function extensionality, proposition extensionality and UIP.
%
Unfortunately, Hofmann's setoid model only interprets a fragment of type theory,
in particular it does not support universes or large elimination (the ability to
define types by pattern-matching on values) which makes it impossible to even
prove that \( 0 \neq 1 \).
%

{\color{red} [Its spiritual successor is observational type theory, which has
been justified both by translations or by logical relations.
%
But OTT is definitionally proof irrelevant, (here I need to talk
about having a specific sort for propositions)
which means that it does not support Acc-elim.
Plus, it does not have unique choice.]}

\paragraph{Contributions}

In this paper, we construct a new universe of setoids which does not require the
equalities to be definitionally proof-irrelevant.
%
We then use our universe of setoids to construct two extensions of Hofmann's
setoid model:
%
a \emph{predicative} model which uses \( \Type \)-valued equalities and can be
formulated in plain Martin-Löf Type Theory (MLTT), and an \emph{impredicative}
model which uses \( \Prop \)-valued equalities and can be formulated in the
Calculus of Inductive Constructions (CIC).
%
The predicative model supports all the type formers of MLTT plus quotient
types, and validates function and proposition extensionality, as well as the
principle of unique choice.
%
The impredicative model adds support for impredicative quantification and
large elimination of the accessibility predicate, but it does not validate the
principle of unique choice.

Because our models are constructed in plain intensional type theory, they
show that extensionality principles and weak choice principles do not add any
consistency strength to MLTT or CIC.
%
By combining this result with the conservativity of extensional type theory over
intensional type theory with extensionality principles~\citep{hofmann1995conservativity, winterhalter_ett},
we obtain that the same is true for the equality reflection rule.
%
More generally, any arithmetical theorem that can be proven in MLTT or CIC
extended with extensionality principles, weak choice principles or the equality
reflection rule can also be derived without them.

%% {\color{red} [Description of the techniques.
%% In doing so, we will methods for reducing sophisticated induction principles such
%% as IR and II to plain inductive families.]}
%% %
%% In doing so, we develop new methods to eliminate the use of induction-recursion
%% and recursion-recursion, which allows our construction to be defined only using
%% indexed inductive families.

\paragraph{Plan of the paper}

{\color{red} The usual.}

\subsection{Related work}

The setoid model was originally conceived by \citet{hofmann95} in order to
extend intensional type theory with extensionality principles and quotient
types.
%
However, this first iteration did not support universes or large elimination.
%
Building on this work, \citet{altenkirch99} used a type theory extended with a sort
of definitionally proof-irrelevant propositions to give a stricter version of
the setoid model, and equipped it with a universe of simple types which contains
only non-dependent function types.
%
This already provides a limited form of large elimination, and in particular this
is enough to prove that \( 0 \neq 1 \).
%
Finally, \citet{setoid_univ} study several methods to construct a universe that
supports dependent type formers, using further extensions of the ambient type
theory such as induction-recursion, induction-induction or a strengthened \( J \)
eliminator for the equality. Compared to their work, our construction only needs
plain indexed inductive types -- in fact, it does not even require definitional
proof-irrelevance.

Meanwhile, \citet{altenkirchAl:plpv2007} explored a different way to mix extensionality and
intensionality with Observational Type Theory (OTT), an intensional type theory
whose equality type is replaced with a definitionally proof-irrelevant setoid
equality.
%
This strategy allows them to implement quotient types and extensionality
principles, while preserving the good computational properties of their system.
%
However, \citet{xtt} observe that the utility of quotient types is substantially
diminished in the absence of unique choice.
%
Subsequently, the community explored several points in the design space of
observational type theories:
%
XTT~\citep{xtt} uses a cubical syntax to implement extensionality principles,
and implements universe induction, but still lacks unique choice.
%
\( \mathsf{CIC}^{\mathrm{obs}} \)~\citep{pujet:hal-03367052} adds an
impredicative sort of propositions and a scheme of indexed inductive definitions
to bring their system closer to CIC, but they are not able to implement large
elimination of the propositional accessibility predicate.
%
Our work is able to handle unique choice for the predicative model, and large
elimination of accessibility for the impredicative model.

%
%% Ultimately, these three systems -- OTT, XTT and \( \mathsf{CIC}^{\mathrm{obs}} \)~--
%% all rely on an extensional theory for their proofs of canonicity, and thus they
%% cannot be said to reduce extensionality principles to an intensional foundation.

%% While the setoid models and the observational type theories both support
%% quotient types, \citet{xtt} observes that their utility is substantially
%% diminished in the absence of unique choice.

\section{Hofmann's setoid model}

\subsection{Our metatheory}

\begin{itemize}
\item The usual CIC with indexed inductive families
\item We will need an impredicative Prop only if we want impredicativity
\end{itemize}

\subsection{Hofmann's model}

\begin{itemize}
\item Blah blah
\end{itemize}

\section{Definition of the universe}

\subsection{Inductive-Recursive intuition}

\begin{itemize}
\item The fundamental idea is that we want to define the universe inductively, while defining the setoid equality by recursion
\item That is technically double small IR
\item in Kaposi et al the suggest to eliminate via induction-induction
\end{itemize}

\subsection{A Plan in Three Steps}

\begin{itemize}
\item Start like you would with small IR
\item Except that here we are trying to do double IR, which is technically not small
\item Therefore, we use this weird encoding that still somehow works
\item Talk about the elimination principle that we get for the universe
\item Probably it has the right computation rules? Maybe look at strong vs weak in Kaposi et al
\end{itemize}

\subsection{Transitivity and Coercion}

\begin{itemize}
\item Mutually define transitivity, coercion, and the equality of coercion
\item Do it forward and backward simultaneously so we can pass Coq's guard condition
\end{itemize}

\subsection{Second universe and embedding}

\begin{itemize}
\item Now that the universe is a proper setoid, we can do a second universe level
\item Put an embedding constructor, rinse and repeat
\item Obviously we can get a full hierarchy
\item No cumulativity though :shrug: maybe we can define emb as a function instead, which might be a tad bit cooler?
\end{itemize}

\section{W Types and Recursion-Recursion}

\subsection{W Types and their equality}

\begin{itemize}
\item We have to restrict a tiny bit the sup to only use extensional functions
\item The equality is inductively generated
\end{itemize}

\subsection{Eliminating W Types}

\begin{itemize}
\item We need to define the eliminator while simultaneously proving it extensional
\item Not so easy actually! The natural way to do it is via double recursion-recursion
\item While simple recursion-recursion is trivial to eliminate (just do a sigma lol),
  double recursion recursion is more technical
\end{itemize}

\subsection{Graph encoding}

\begin{itemize}
\item Define the graph inductively
\item Show that it is extensional
\item Then show that you can always find
\end{itemize}

\section{Quotients and Unique Choice}

\subsection{Adding Quotient Types}

\begin{itemize}
\item Not very difficult tbh, just add them
\end{itemize}

\subsection{The Principle of Unique Choice}

\begin{itemize}
\item Allows for effective quotients
\item Logically very powerful, very close to replacement
\item Impossible for proof-irrelevant OTT, since there's no info in equality
\item Unfortunately the naive version causes some universe issues
\item Maybe we can use Pédrot's trick
\item But for that we would need a proper normalisation proof for an impredicative theory
\end{itemize}

\subsection{A Predicative Theory}

\begin{itemize}
\item Instead, we propose a new model with a predicative hierarchy of propositions
\item That way, we can get a principle of unique choice which actually works
\item In this model, functions and functional graphs are exactly the same
\end{itemize}

\section{Impredicative Quantification and Accessibility}

\subsection{Flirting with Russell's paradox}

\begin{itemize}
\item recall why impredicativity is nice?
\end{itemize}

\subsection{Well-founded induction and Weak Choice Principles}

\begin{itemize}
\item Explain how acc elimination works
\item Allowed in Lean with strict proof irr, failure of normalisation
\item Disallowed in OTT, which implies that we cannot extract a function that normalises System F
\item It implies some small bits of choice!
\end{itemize}

\subsection{Adding Acc elimination}

\begin{itemize}
\item We need to do the graph trick again
\item but it works, nice!
\end{itemize}

\section{Conclusion and Future Work}

\begin{itemize}
\item Recall the results: we only need inductive family for our model of extensionality!
\item Maybe some applications stolen from Hofmann's thesis?
\item Maybe make it into a proper type theory, OTT style
\item We are not going to have cast on refl, though :shrug:
\item Normalisation proofs, normalisation proofs, normalisation proofs
\item Making unique choice work with impredicativity?
\end{itemize}

%\newpage
\bibliography{biblio}

\end{document}
